\chapter{Introdução}\label{chp:INTRODUCAO}

Na era da informa\c{c}\~{a}o, a dissemina\c{c}\~{a}o de novos meios de comunica\c{c}\~{a}o dentro do ambiente universit\'{a}rio, com o intuito de facilitar a comunica\c{c}\~{a}o entre discentes e docentes, vem se tornando cada vez mais relevante. Nesse cen\'{a}rio, os agentes conversacionais podem ser ferramentas importantes e \'{o}timos intermediadores para a busca de informa\c{c}\~{o}es de \^{a}mbito acad\^{e}mico, como respostas a FAQs relacionadas a matr\'{\i}culas, grade curricular e eventos, facilitando e otimizando o tempo das secretarias, que outrora possu\'{\i}am o papel de entregar esse tipo de resposta.

Conforme j\'{a} observado em estudos acad\^{e}micos, tal como (i) a monografia \textit{``Solu\c{c}\~{a}o Chatbot no Ambiente Acad\^{e}mico da UFRJ''} \cite{ufrj}, onde \'{e} apresentada uma solu\c{c}\~{a}o baseada em regras dentro do ambiente acad\^{e}mico da UFRJ, focada na entrega de informa\c{c}\~{o}es sobre matr\'{\i}cula, calend\'{a}rio acad\^{e}mico, atestados e declara\c{c}\~{o}es, e tamb\'{e}m (ii) o trabalho \textit{``Um Chatbot Especializado para o Contexto da Universidade Federal de Ouro Preto''} \cite{ufop}, com aplica\c{c}\~{o}es similares e tecnologias mais recentes, observa-se que a utiliza\c{c}\~{a}o de chatbots pode ser aplicada com esses fins.

Atualmente, na Universidade Federal Rural do Rio de Janeiro, as atas de reuni\~{o}es acad\^{e}micas do Departamento de Ci\^{e}ncia da Computa\c{c}\~{a}o armazenam informa\c{c}\~{o}es essenciais sobre o desenvolvimento do curso e decis\~{o}es acad\^{e}micas que afetam os alunos e os professores, representando um hist\'{o}rico valioso para preserva\c{c}\~{a}o e consulta.

Nesse contexto, foi observada a necessidade de uma ferramenta que facilite a busca por atas — atas que contenham informa\c{c}\~{o}es pertinentes sobre assuntos espec\'{\i}ficos ou buscas mais gen\'{e}ricas. Com fins de pesquisa e levantamento do hist\'{o}rico de reuni\~{o}es dentro do departamento e considerando que a quantidade de atas cresce continuamente, realizar a busca manual torna-se cada vez mais dif\'{\i}cil. Al\'{e}m disso, a aplica\c{c}\~{a}o de um sistema de consulta de atas demanda manuten\c{c}\~{a}o, aprendizado por parte do operador e implanta\c{c}\~{a}o adequada ao ambiente acad\^{e}mico.

\section{Objetivo}

Diante disso, o problema central deste trabalho \'e permitir que perguntas realizadas por discentes, docentes ou o p\'{u}blico acad\^{e}mico no geral sejam interpretadas e mapeadas para as atas que mais se adequam ao conte\'{u}do de suas buscas, com o interm\'{e}dio de um chatbot de regra para a comunica\c{c}\~{a}o e um sistema de busca dentro de uma base de dados composta por essas atas.

Este artigo visa (i) investigar os principais tipos de chatbot dispon\'{\i}veis para a implanta\c{c}\~{a}o, (ii) analisar as tecnologias envolvidas em sua constru\c{c}\~{a}o e desenvolvimento e (iii) realizar um estudo mais aprofundado sobre chatbots baseados em regras, considerando seus principais benef\'{\i}cios e aplica\c{c}\~{o}es e a usabilidade dentro do problema apresentado.

\section{Defini\c{c}\~{a}o e Origem dos Chatbots}

Segundo a IBM \cite{IBMChatbots}, chatbot é um programa de computador desenvolvido para simular a conversação humana com o usuário final. Embora nem todos os chatbots sejam equipados com inteligência artificial, os modelos mais recentes utilizam técnicas avançadas, como o processamento de linguagem natural, para compreender as questões apresentadas pelos usuários e automatizar as respostas.


A história dos chatbots teve início na década de 1960, através do desenvolvimento do programa ELIZA, idealizado por Joseph Weizenbaum no Massachusetts Institute of Technology (MIT). ELIZA foi um chatbot projetado para simular uma conversa com um psicoterapeuta, funcionando com base em regras simples de construção de palavras e scripts pré-definidos \cite{weizenbaum1966eliza}. ELIZA não compreendia de fato o conteúdo das mensagens, mas conseguia simular interações humanas de maneira revolucionária para a época, sendo sua criação um marco do início da exploração de agentes conversacionais no campo da inteligência artificial (IA) e da linguística computacional.

Outro marco na evolução dos chatbots ocorreu em 1972, quando o psiquiatra Kenneth Colby, na universidade de Stanford desenvolveu o PARRY, modelo que, diferentemente do ELIZA, simulava o comportamento de um paciente com esquizofrenia paranoide. O PARRY foi testado em uma versão do Teste de Turing \cite{turing1950} (Onde Turing propõe um teste para avaliar se uma máquina pode exibir comportamento inteligente indistinguível do humano, por meio do que chamou de "Imitation Game", também conhecido como Teste de Turing), onde foi submetido a julgamento por psiquiatras que interagiram com o programa, analisando suas transcrições. Os especialistas identificaram corretamente cerca de 48\% das trocas, valor que demonstra que, mesmo para profissionais da área, o PARRY se passava, na maioria das situações, por um paciente humano \cite{colby1972parry}.

Essa capacidade de simulação deixa evidente o avanço  no uso de agentes conversacionais como ferramentas de investigação psicológica, mesmo em um tempo onde a tecnologia não era tão avançada como hoje, reforçando a relevância do PARRY como uma etapa importante na história da IA conversacional.

\section{Evolução dos Chatbots}

Após os primeiros experimento com o ELIZA e o PARRY, o desenvolvimento de chatsbots avançou consideravelmente com a criação do Artificial Linguisticc Internet Computer Entity (ou simplesmente A.L.I.C.E), em 1995, por Richard Wallace. Diferente de seus predecessores, A.L.I.C.E foi projetado com uso de uma linguagem específica chamada AIML (Artificial Intelligente Markup Language), que permitiu a construção de padrões de entrada e respostas de forma modular e escalável, facilitando sua expansão e adaptação para diferentes tipos de diálogos \cite{wallace2009alice}.

Mesmo que A.L.I.C.E não utilizasse técnicas de aprendizado de máquinas, sua estrutura baseada em regras, somado ao grande volume de dados escritos manualmente, possibilitava interações elaboradas, tendo A.L.I.C.E conquistando reconhecimento internacional após vencer 3 vezes o Prêmio Loebner (2000, 2001 e 2004), competição baseada no Teste de Turing \cite{turing1950} que avalia o grau de humanização de um chatbot \cite{shum2018eliza}.

A partir de 2011, com o lançamento da Siri pela Apple, integrada ao iPhone 4S, iniciou-se uma nova era no desenvolvimento de chatbots inteligentes, marcando a popularização dos Assistentes Pessoais Inteligentes (APIs) no cotidiano dos usuários. A Siri foi pioneira ao trazer uma interface de voz capaz de compreender comandos naturais e executar tarefas diversas, integrando múltiplas fontes de dados e sensores embarcados no dispositivo \cite{hoy2018}.

Ao longo da década surgiram diversos Assistentes Pessoais Inteligentes, como a Cortana da Microsoft (2014), o Google Assistente (2016) e a Alexa da Amazon (2014). Esses agentes são capazes de integrar informações provenientes de múltiplos sensores, como localização, clima, movimento, toques e gestos, além de acessarem diversas fontes de dados pessoais, como músicas, calendários, e-mails e perfis sociais \cite{shum2018eliza}.

Nos últimos anos alcançamos um novo patamar em relação à evolução dos Chatbots, sendo adotados modelos de linguagem de grande escala, conhecidos como LLMs (Large Language Models) \cite{IBMLLM}. Esses modelos utilizam técnicas de Deep Learning para compreender e gerar informações de maneira natural e respeitando contextos específicos da conversa.

Ao contrário dos chatbots mais tradicionais, que funcionam com base em fluxos fixos ou regras pré-programadas, os modelos de linguagem de larga escala (LLMs) são treinados com enormes volumes de texto e conseguem lidar com tarefas bem mais complexas. Eles podem responder perguntas abertas, escrever textos longos e bem estruturados, resumir documentos, traduzir idiomas, explicar conceitos técnicos e até simular estilos de fala ou perfis de personalidade. Essa flexibilidade e capacidade de adaptação a diferentes situações tornam esses modelos especialmente úteis em áreas como educação, atendimento ao cliente, saúde e produtividade no ambiente corporativo \cite{bommasani2021opportunities}.

\section{Tipos de Chatbots}

A IBM, através da publicação "Seis tipos de chatbots e como escolher o ideal para sua empresa" \cite{IBMChatbots} nos mostra alguns tipos de chatbots disponíveis no mercado hoje, sendo eles:

\subsection{Chatbots Baseados em Menus ou Botões}
Esses chatbots compõem o nível mais básico de chatbots. A interação com o usuário se dá com a seleção de um script que se encaixa melhor nas necessidades atuais do utilizador, podendo haver mais opções conforme as escolhas forem sendo feitas, até chegar na opção mais adequada. Esses chatbots operam como árvores de decisões e são bons para tarefas transacionais. São focados na resolução de funcionalidades simples e extremamente úteis para a resposta de perguntas repetitivas, previsíveis e diretas. A dificuldade é observada com solicitações menos genéricas, pois possuem opções de respostas pré-definidas e limitadas. Além disso, caso a necessidade do usuário não esteja disponível nas opções do menu, o chatbot se torna inútil.

\subsection{Chatbots Baseados em Regras}
Possuem a funcionalidade da árvore de decisão, assim como os chatbots baseados em menu, adicionando também regras de lógica condicional para o desenvolvimento de fluxos de automação de conversa. Atuam essencialmente com FAQs interativas, necessitando de uma quantidade de combinações pré-definidas de perguntas e respostas para a interação com o usuário.
Operam com detecção básica de palavras-chave, sendo fáceis de treinar e funcionando bem também quando recebem perguntas em um escopo pré-definido, tendo dificuldade para lidar com perguntas mais complexas e que não foram previstas no escopo.

\subsection{Chatbots Impulsionados por IA}
Possuem capacidade de compreensão das perguntas do usuário com recursos de IA e Entendimento de Linguagem Natural (NLU), conseguindo detectar rapidamente as informações relevantes a partir das interações com o usuário, tornando a interação mais conversacional.
Possuem a possibilidade de buscar mais informações com o usuário para o entendimento correto do escopo fornecido e solicitado. Também podem mostrar opções de respostas onde o usuário pode escolher a que mais se encaixa no contexto encontrado.
Utilizam algoritmos de aprendizado de máquina em sua implementação, permitindo entendimento, aprendizagem e desenvolvimento.

\subsection{Voice Chatbots}
São ferramentas de conversação que permitem a interação falada em alternativa à interação textual, podendo oferecer as mesmas funcionalidades que os chatbots de IA convencionais, mas são implementados com tecnologia Text-to-Speech e Speech-to-Text.
Com a aplicação do Processamento de Linguagem Natural e interações com tecnologias de computação e telefonia, os chatbots por voz conseguem compreender perguntas faladas, realizar a análise dos usuários e apresentar as respostas em um tom de conversa, aumentando o engajamento e a imersão do usuário.

\subsection{Chatbots de IA Generativa}
Podem oferecer funcionalidades mais avançadas que os anteriores, possuindo maior fluência na compreensão da linguagem comum e maior adaptação à linguagem do usuário, podendo gerar conteúdo novo como resultado. Esse conteúdo pode ser texto, imagens ou sons de alta qualidade, baseados em Large Language Models (LLMs) nos quais são treinados.
Apresentam maior personalização aos usuários, aumentando o alcance e a velocidade no atendimento. As interfaces de chatbots com IA generativa podem reconhecer, resumir, traduzir, prever e criar conteúdo em resposta à consulta de um usuário.

\subsection{Chatbots Híbridos}
Chatbots híbridos são sistemas de IA conversacional que combinam lógica baseada em regras e recursos baseados em aprendizado de máquina, proporcionando uma experiência mais versátil ao lidar com tarefas de diferentes níveis de complexidade. Funcionam com um conjunto de regras e scripts pré-definidos, apresentando estrutura, enquanto a IA possui um potencial de aprendizado para interações mais complexas.

\section{Escolha do modelo de Chatbot}
Diante da necessidade de uma ferramenta que faça a interação com o usuário e realize a busca da ata mais adequada de acordo com o questionamento, propõe-se a utilização de um sistema que se baseia em consultas a uma base de dados estruturada, que irá se compor pelas atas, e que será realizada por meio de uma interface conversacional integrada a um canal de comunicação utilizado no ambiente universitário. Optou-se pela aplicação de um chatbot baseado em regras, pela capacidade do estabelecimento de um escopo e uma estrutura de diálogo pré-definida, permitindo o mapeamento de respostas prontas com base em fluxos lógicos bem definidos.
Para viabilizar a recuperação eficiente dos documentos mais adequados com base nas consultas feitas pelo usuário, será implementado um mecanismo de busca por similaridade semântica utilizando o Elasticsearch \cite{IBMElasticsearch}. Essa tecnologia permite a indexação de grandes volumes de texto e a realização de buscas avançadas com alto desempenho, o que é fundamental para garantir a precisão e a velocidade no retorno das informações. O Elasticsearch será configurado para processar e ranquear as atas com base na similaridade entre o conteúdo dos documentos e a pergunta inserida, aumentando significativamente a relevância das respostas apresentadas pelo sistema.

\section{Conceitos Relavantes}

\subsection{Large Language Models (LLMs)}
Segundo a IBM \cite{IBMLLM}, LLM (Large Language Model) é um modelo de Machine Learning aplicado na Inteligência Artificial generativa, capaz de interpretar diferentes funções. Os LLMs são treinados com imensas quantidades de dados para o reconhecimento e geração de imagens, textos e conteúdos diversos.
Tais bases de dados são fundamentais para impulsionar casos de uso e aplicações. Foram projetados para entender e gerar texto como um humano, com base na vasta quantidade de dados utilizados. São capazes de inferir a partir do contexto, gerar respostas coerentes e contextualmente relevantes, realizar traduções para outros idiomas, resumos de textos, respostas a perguntas e até auxiliar em tarefas criativas.
Possuem parâmetros que permitem a captura de padrões complexos na linguagem e representam um avanço considerável no processamento de linguagem natural e na inteligência artificial. Há exemplos amplamente disseminados como o ChatGPT da OpenAI, o LLaMA da Meta, o Watson da IBM e o Gemini do Google.

Os LLMs são constituídos por várias camadas de redes neurais, com parâmetros que podem ser ajustados durante o treinamento e que são aprimorados ainda mais por uma camada conhecida como mecanismo de atenção, que se concentra em partes específicas dos conjuntos de dados.
Geralmente, são baseados em uma arquitetura de transformer, como o Transformador Generativo Pré-Treinado (GPT), projetado para lidar com dados sequenciais, como entradas de texto.

Durante o processo de treinamento, os modelos realizam ações de previsão com base em uma pontuação de probabilidade de recorrência de palavras que foram tokenizadas (divididas em sequências menores de caracteres). Esses tokens são transformados em embeddings (representações numéricas do contexto), que são utilizados para gerar saídas contextualizadas.
\subsection{Elasticsearch}
Segundo a IBM \cite{IBMElasticsearch}, Elasticsearch é um motor de busca e análise de código aberto, baseado na biblioteca Apache Lucene. Oferece recursos de busca escaláveis, podendo armazenar dados estruturados em formatos especializados para buscas otimizadas. Possui também um design de API RESTful, o que proporciona flexibilidade para a análise de dados e facilidade para a integração com aplicações. Trata-se de uma solução de busca textual rápida e amplamente disseminada, com escalabilidade horizontal e compatibilidade com diversas linguagens de programação, incluindo Java, Python, .NET, PHP, entre outras, Usa o algoritmo BM25 para ranquear resultados com base na relevância da correspondência textual e possui suporte a linguagem natural e análise linguística, com ferramentas internas de análise morfológicas, permitindo a criação de analisadores personalidades afim de aumentar a precisão da busca textual.


\section{Botpress}

Foi realizada uma análise dos principais tipos de Chatbots baseados em regras (ou chatbots mais inteligentes, que poderiam ser utilizados no contexto baseado em regras), tendo sido optado pela utilização do \textbf{Botpress}, principalmente por não ser necessário nenhum tipo de licenciamento. Também foi cogitado a utilização de versões grátis de modelos pagos, porém o limite de conversação inviabilizaria a continuidade do projeto.

O Botpress é uma ferramenta de código aberto usada para criar chatbots inteligentes. Para esse projeto, foi escolhida a versão v12, principalmente por ser totalmente gratuita e permitir que o sistema seja hospedado localmente, sem depender de servidores externos ou planos pagos \cite{botpress-about}. Isso garante mais autonomia, controle sobre os dados e liberdade para personalizar conforme as necessidades do trabalho.

A plataforma oferece uma interface visual, \textit{drag-and-drop}, para montar os fluxos de conversa, além de permitir integrações com APIs, uso de processamento de linguagem natural (NLP) e armazenamento local de informações \cite{botpress-docs}. Essas funcionalidades tornam o Botpress uma boa opção para contextos educacionais e acadêmicos, justificando a escolha.

\section{Organização do Trabalho}

Nesta parte, será explicada como será a organização do trabalho.